\section{Implementation}
\label{sec:beer-implementation}

% Plan: chapters/beer/plan.md §7.6. The single, explicitly subordinate implementation section.
% Migrated from the former lean-implementation-architecture.tex and binders-names-contexts.tex.

The mathematical rules of \Cref{sec:beer-encoding-rules} and the soundness statements of
\Cref{sec:beer-soundness} are realized in Lean as a state-passing computation. This section records
how, and---more to the point---why the state never surfaces in the mathematical development.

\subsection{The two-pass pipeline and the encoder monad}
\label{sec:lean-implementation-architecture}

\beer{} is structured as a two-pass pipeline that threads a single mutable state
functionally, through a \emph{monad}. Recall that a monad organizes effectful
computations: it fixes a notion of effect---here, reading and updating a state and
failing with an error---together with the sequencing discipline that Lean's
\leaninl{do}-notation renders implicit, where \leaninl{let x ← c} runs the
computation \leaninl{c}, propagates its effects, and binds its result to
\leaninl{x}. Each of the two passes runs in one such monad.

The first pass, the \emph{decoder}, parses an Atelier~B proof-obligation file
into the internal concrete representation of \Cref{ch:b}. It
runs in the monad {\small\leancodeptr{Decoder}{POGReader/Basic.html\#Decoder}}, defined as
\leaninl{StateT DecoderState (ExceptT String IO)}: a state carrying the B
environment under construction, over exceptions, over \leaninl{IO} for reading the
file.

The second pass, the \emph{encoder}, is the subject of this thesis. It transforms
the decoded B environment into an SMT-LIB script and runs in
{\small\leancodeptr{Encoder}{Encoder/Basic.html\#Encoder}}, defined as
\leaninl{StateT EncoderState (Except String)}. Its state records everything the
translation must track:
\begin{leanlisting}
structure EncoderState where
  env   : SMT.Env          -- fresh-variable counter, used names, emitted script
  types : SMT.TypeContext  -- the running SMT typing context Γ
\end{leanlisting}
namely the running SMT typing context~$\Gamma$ (the field \leaninl{types}), a
monotone fresh-variable counter together with the set of names already in use, and
the SMT declarations and assertions accumulated so far---the eventual output.
The encoder's effects are exactly these: consulting and
extending~$\Gamma$, allocating fresh names, emitting declarations and assertions,
and raising an error on an ill-typed or unsupported input.

Crucially, the encoder is \emph{pure}: unlike the decoder, its stack contains no
\leaninl{IO}. It therefore denotes a total function from a B environment to either
an error or an SMT script paired with a final state. This is what licenses the
mathematical style adopted throughout this chapter: the correctness
statements of \Cref{sec:beer-soundness} are equalities between the
denotation of the emitted term and that of the source term, in which the state
threaded through every intermediate step never appears. In the Lean development the
state-passing obligations are discharged separately, by reasoning about the monadic
primitives, \eg \leaninl{get}, \leaninl{modifyGet}, and \leaninl{throw}, in a
Floyd--Hoare-style~\cite{floyd1967meanings,hoare1969axiomatic} weakest-precondition
calculus~\cite{dijkstra1975guarded} for monadic programs, so that the encoding rules may
be read---and verified---as if they manipulated~$\Gamma$ and the fresh-name supply directly.

\subsection{Handling of binders, names, and contexts}
\label{sec:binders-names-contexts}

Encoding a quantified or functional B construct forces the encoder to invent SMT
names for the variables it binds and to record their types. Both tasks are realized
as effects of the {\small\leancodeptr{Encoder}{Encoder/Basic.html\#Encoder}} monad
(\Cref{sec:lean-implementation-architecture}), acting on two components of its state
that are given deliberately different lifetimes.

\paragraph{Fresh names and the typing context.}
The operation {\small\leancodeptr{freshVar}{Encoder/Basic.html\#SMT.freshVar}}, given a type
$\tau$, returns a name distinct from every name used so far, records it among the
used names, and inserts the binding into the running context~$\Gamma$;
{\small\leancodeptr{addToContext}{Encoder/Basic.html\#SMT.addToContext}} and
{\small\leancodeptr{eraseFromContext}{Encoder/Basic.html\#SMT.eraseFromContext}} insert and
remove bindings explicitly. The fresh-variable counter and the set of used names are
\emph{monotone}---never rolled back---which guarantees that every generated name is
globally unique across the whole run. The typing context~$\Gamma$, by contrast, is
\emph{scoped}: a bound variable is meaningful only while its body is being encoded,
so~$\Gamma$ is extended on entering a binder and restored on leaving it.

\paragraph{Erasing binder variables.}
The interplay is visible in
{\small\leancodeptr{defaultSpecM}{Encoder/Loosening/Rules.html\#defaultSpecM}}, which builds
the predicate expressing that a term takes the canonical default value of its
type. At a function type $\alpha \toS \beta$ the default is taken
pointwise, $\forallS x : \alpha.\ \mathsf{dflt}_{\beta}\bigl(t(x)\bigr)$, which in the
encoder reads:
\begin{leanlisting}
| .fun α β, t => do
    let x    ← freshVar α                     -- inserts x : α into Γ
    let body ← defaultSpecM β (.app t (.var x))
    eraseFromContext x                        -- x occurs only under the ∀ below
    pure (.forall [x] [α] body)
\end{leanlisting}
The fresh $x$ is drawn solely to serve as the universal binder. Since
\leaninl{freshVar} has inserted $x : \alpha$ into~$\Gamma$, but $x$ occurs only under
the $\forallS$ about to be built, it must be erased from~$\Gamma$ before the
predicate is returned: otherwise the result would carry a spurious free constant~$x$
and fail to typecheck in the ambient context, where $x$ is a variable bound by the
quantifier rather than a global symbol.

The same discipline recurs throughout the encoder. Every case of
{\small\leancodeptr{loosenAux}{Encoder/Loosening/Rules.html\#loosenAux_prf}} erases the
$\existsS$- and $\lamS$-binders it introduces once their specification is assembled,
and the universal case of {\small\leancodeptr{encodeTerm}{Encoder/Encoder.html\#encodeTerm}}
additionally snapshots~$\Gamma$ and the emitted declarations before encoding the
body: helper declarations produced while translating under the binder are then
re-scoped \emph{inside} the quantifier---as guarded universals---rather than left as
global constants that would escape their intended scope.

\begin{note}
  Keeping the fresh-name counter monotone while rolling back~$\Gamma$ is a matter of
  correctness, not convenience. Restoring the counter on leaving a binder could
  re-issue a name already attached to an emitted declaration, conflating two
  distinct SMT symbols; keeping it monotone preserves the global freshness invariant
  on which the correctness proofs rely, at the negligible cost of a few unused
  indices.
\end{note}
